\origpage{331}{N. Robertson, P.D. Seymour / Journal of Combinatorial Theory, Series B 90 (2004) 325--385}
\noindent $|\tilde{e}| = 3$, choose $c\in E(\Gamma_i)$ with $e = s_i(c)$, and let $\phi'_i(e) = (\phi_i(c),3)$. If $e\in E_i^2$ and $e$ is internal, let $\phi'_i(e) = (\phi_i(e),2)$. If $e$ is a border edge of $\Gamma'_i$ (and therefore $e\in E_i^2$), but not a short side, let $\phi'_i(e) = (\phi_i(e),\pm 1)$, where $-1$ is chosen if $\gamma'_i(e)\ne \gamma_i(e)$. If $|\tilde{e}| = 2$ and $e\notin E_i^2$, let $\phi'_i(e) = (x,0)$ where $x\in E(\Omega)$ is arbitrary. For each short side $e$ of $\Phi$, let $\phi'_i(e) = e$. It follows easily that for each $i\ge 1$, $(\Gamma'_i,\phi'_i)$ is a $\chi$-coloured painting.

From 3.1, there exists $j>i\ge 1$ and a linear inflation $\sigma'$ of $(\Gamma'_i,\phi'_i)$ in $(\Gamma'_j,\phi'_j)$. We claim that there is a linear inflation $\sigma$ of $\Gamma_i$ in $\Gamma_j$ such that $\phi_i\le \phi_j\circ \sigma$. We construct $\sigma$ as follows. For each $v\in V(\Gamma_i)$ let $\sigma(v) = \sigma'(v)\cap sk(\Gamma_j)$. Since each $\sigma'(v)$ is induced and every edge of $sk(\Gamma'_j)$ either is an edge of $sk(\Gamma_j)$ or is parallel to an edge of $sk(\Gamma_j)$, it follows that each $\sigma(v)$ is a non-null connected induced subgraph of $sk(\Gamma_j)$. For $e\in E_i^3$, let $\sigma(e) = f$, where $\sigma'(s_i(e)) = s_j(f)$ (such an $f$ exists since $\sigma'(s_i(e))$ has the same size as $s_i(e)$). For $e\in E_i^2$, let $\sigma(e) = \sigma'(e)$. Let us verify that $\sigma$ is a linear inflation $\sigma$ of $\Gamma_i$ in $\Gamma_j$ such that $\phi_i\le \phi_j\circ \sigma$. First, by examining the second terms of $\phi'$, it follows that $\sigma(e) = e$ for every short side $e$. Moreover, if $e\in E_i^2$ is not a short side then since the second term of $\phi'_i(e)$ is non-zero, so is the second term of $\phi'_j(\sigma'(e))$, since $\phi'_i(e)\le \phi'_j(\sigma'(e))$, and hence $\sigma'(e)\in E_j^2$. By examining the first term of $\phi'_i(e)$ and using (1), we see that $\phi_i\le \phi_j\circ \sigma$. To complete the proof, we must check that

\hspace*{12pt}(2) \textit{For each $e\in E(\Gamma_i)$ and $1\le k\le |\tilde{e}|$, if $v$ is the $k$th end of $e$ then $\sigma(v)$ contains the $k$th end of $\sigma(e)$.}

\textit{Subproof.} Now $v$ is the $k$th term of $\gamma_i(e)$. If $|\tilde{e}|=3$, then $\gamma'_i(s_i(e)) = \gamma_i(e)$, and therefore $V(\sigma'(v)) = V(\sigma(v))$ contains the $k$th term of $\gamma'_j(\sigma'(s_i(e))) = \gamma_j(\sigma(e))$ as required. Similarly, if $|\tilde{e}|=2$ and $e$ is internal or a short side of $\Phi$, then $\gamma'_i(e) = \gamma_i(e)$, and therefore $V(\sigma'(v)) = V(\sigma(v))$ contains the $k$th term of $\gamma'_j(\sigma'(e)) = \gamma_j(\sigma(e))$ as required. Finally, suppose that $e$ is a border edge of $\Gamma_i$ and not a short side of $\Phi$. Let $S$ be the long side of $\Phi$ with $e,\sigma(e)\subseteq S$. Since the second terms of $\phi'_i(e)$ and $\phi'_j(\sigma(e))$ are equal, it follows that $\gamma'_i(e) = \gamma_i(e)$ if and only if $\gamma'_j(\sigma(e)) = \gamma_j(\sigma(e))$, and so there exists $l$ ($l=1$ or $2$) such that the $k$th term of $\gamma_i(e)$ is the $l$th term of $\gamma'_i(e)$, and the $k$th term of $\gamma_j(\sigma(e))$ is the $l$th term of $\gamma'_j(\sigma(e))$. Therefore $v$ is the $l$th term of $\gamma'_i(e)$, and so $\sigma'(v)$ contains the $l$th term of $\gamma'_j(\sigma(e))$. Consequently $\sigma(v)$ contains the $k$th term of $\gamma_j(\sigma(e))$. This proves (2).

From (2), it follows that $\sigma$ is an inflation of $\Gamma_i$ in $\Gamma_j$. Moreover, it is \del{clearly} linear, and satisfies $\phi_i\le \phi_j\circ \sigma$, as required. \fix{We check the three conditions in the definition of ``linear''. For the first, let $e \in E(\Gamma_i)$ and let $\Theta$ be a cuff. If $|\tilde{e}| = 3$, then $e$ and $\sigma(e)$ are internal, since border edges have size 2, and so neither borders $\Theta$. If $|\tilde{e}| = 2$, then $e$ and $\sigma(e) = \sigma'(e)$ are edges of $\Gamma'_i$ and $\Gamma'_j$, and $e$ borders $\Theta$ if and only if $\sigma'(e)$ does, since $\sigma'$ is linear. For the second, let $e'$ be a border edge of $\Gamma_j$ with $e' \notin \sigma(E(\Gamma_i))$. Then $e'$ is a border edge of $\Gamma'_j$, and it is not a short side, since $\sigma(s) = s$ for every short side $s$; so the second term of $\phi'_j(e')$ is $\pm 1$. We claim that $e' \notin \sigma'(E(\Gamma'_i))$. For suppose that $e' = \sigma'(c)$. Then $c$ has size 2, since $e'$ has. If $c \in E_i^2$, then $\sigma(c) = \sigma'(c) = e'$, which is impossible. Otherwise $c$ is an edge of $sk(\Gamma_i)$ not in $E(\Gamma_i)$, and so the second term of $\phi'_i(c)$ is $0$; but $\phi'_i(c) \le \phi'_j(e')$, and $R$ is ordered by equality, a contradiction. This proves the claim. Since $\sigma'$ is linear, it follows that $e'$ is an edge of $\sigma'(v)$ for some $v \in V(\Gamma_i)$. Both ends of $e'$ are in $V(\sigma'(v)) = V(\sigma(v))$, and $e'$ is an edge of $sk(\Gamma_j)$; so $e'$ is an edge of $\sigma(v)$, since $\sigma(v)$ is an induced subgraph of $sk(\Gamma_j)$. For the third, let $e$ border a cuff $\Theta$. Then $e \in E_i^2$ and $\sigma(e) = \sigma'(e)$. As in the proof of (2), either $\gamma'_i(e) = \gamma_i(e)$ and $\gamma'_j(\sigma(e)) = \gamma_j(\sigma(e))$, or $\gamma'_i(e)$ and $\gamma'_j(\sigma(e))$ are the reverses of $\gamma_i(e)$ and $\gamma_j(\sigma(e))$, since these marches have two terms. In the first case the condition for $\sigma$ is the condition for $\sigma'$. In the second case, orient $\Theta$ so that the tail of $e$ in $\Gamma_i$ immediately precedes $e$. Under the reverse orientation, the tail of $e$ in $\Gamma'_i$ immediately precedes $e$; so, since $\sigma'$ is linear, the tail of $\sigma(e)$ in $\Gamma'_j$ immediately precedes $\sigma(e)$ under the reverse orientation. Hence, under the first orientation, the tail of $\sigma(e)$ in $\Gamma_j$ immediately precedes $\sigma(e)$.}\fixnote{fix:2.1-linear-proof-added}{The paper says only that $\sigma$ is ``clearly'' linear. The argument above is the one used in the formalization.} $\square$

\sectionhead{4. The induction}

The remainder of the paper is devoted to proving 3.1. We proceed by an induction on $\chi$ which we shall explain in this section.
